\documentclass[leqno]{article}
%%
%% test leqno
%%
\usepackage{amssymb}
\usepackage{keytheorems-thmmarks}

\newkeytheoremstyle{heartstyle}{
    qed=\ensuremath{\heartsuit}
    }
\newkeytheoremstyle{diamondstyle}{
    qed=\ensuremath{\diamond},
    bodyfont=\normalfont
    }
\newkeytheoremstyle{bulletstyle}{
    qed=\ensuremath{\bullet},
    bodyfont=\normalfont
    }

\newkeytheorem{theorem}[
    style=heartstyle,
    parent=section,
    ]
\newkeytheorem{lemma}[ % not endmarked
    parent=section
    ]
\newkeytheorem{remark}[
    style=diamondstyle,
    parent=section,
    ]
\newkeytheorem{note}[
    style=bulletstyle,
    parent=section,
    ]



\begin{document}

\section{Basic tests}
\begin{theorem}
  Case 1: Environment ends with simple text.
\end{theorem}
\begin{theorem}
  Case 2: Environment ends with display
  \[
    1 + 1 = 2.
  \]
\end{theorem}
\begin{theorem}
  Case 3: Environment ends with equationlike (numbered)
  \begin{equation}
    1 + 1 = 2.
  \end{equation}
\end{theorem}
\begin{theorem}
  Case 4: Environment ends with
  \begin{itemize}
  \item a
  \item list.
  \end{itemize}
\end{theorem}



\section{Other simple tests}

These test multiple or nested boxes inside of theoremlike environments.

\begin{remark}
  Case 1: Environment ends with simple text.
  But first there is an equation
  \begin{equation}
    1+1 = 2
  \end{equation}
  and then the simple text.
\end{remark}
\begin{remark}
  Case 2: Environment ends with display
  \begin{enumerate}
  \item but first
  \item a list
  \end{enumerate}
  \[
    1 + 1 = 2.
  \]
\end{remark}
\begin{remark}
  Case 3: Environment ends with align
  following
  \begin{itemize}
  \item a
  \item list
  \end{itemize}
  and
  \[
    \texttt{ display math }
  \]
  \[
    \texttt{ again }
  \]
  \begin{align}
    x & = 2\\
    y & = 3. \quad \text{this is a very wide box that uses many words to fill the line}
  \end{align}
\end{remark}

\begin{remark}
  Case 4: Environment ends with
  \begin{itemize}
  \item a
  \item list.
  \item that is nested
    \begin{enumerate}
    \item inside another
    \item list
    \end{enumerate}
  \end{itemize}
\end{remark}

\begin{note}
  Case 4: Environment ends with
  \begin{itemize}
  \item a
  \item list.
    \begin{enumerate}
    \item and there is
    \item a nested list
    \end{enumerate}
  \item but the outer list has
  \item the final item 
  \end{itemize}
\end{note}

\section{some other messy tests}
\begin{theorem}
regular text
\end{theorem}

\begin{theorem}
  This ends in a display
  \[
    3^2 + 4^2 = 5^2.
  \]
\end{theorem}

\begin{theorem}
  This ends in a display
  \[
    3^2 + 4^2 = 5^2
  \]
  more text
  \begin{align}
    a&=b \\
    c&=d
  \end{align}
%\[a=b\]
\end{theorem}

\begin{theorem}
  This ends in a display
  \[
    3^2 + 4^2 = 5^2.
  \]
more text
\begin{align*}
a=b
\end{align*}
%\[a=b\]
abc
\end{theorem}

\begin{theorem}
  This ends with
  \begin{enumerate}
  \item an
  \item enumerated
  \item list.
  \end{enumerate}
\end{theorem}

\begin{theorem}
  This ends with
  \begin{enumerate}
  \item an
  \item enumerated
  \item list
  \item ending in equation:
  \begin{align*}
  a&=b
  \end{align*}
  \end{enumerate}
\end{theorem}

\begin{theorem}
  This ends with
  \begin{enumerate}
  \item an
  \item enumerated
  \item list.
  \begin{enumerate}
  \item abc
  \end{enumerate}
  \item blub
  \end{enumerate}
\end{theorem}

\begin{theorem}
  This ends with
  \begin{enumerate}
  \item an
  \item enumerated
  \item list.
  \begin{enumerate}
  \item abc
  \end{enumerate}
  \end{enumerate}
\end{theorem}

\subsection{tests of nested theoremlike environments}

\begin{theorem}
regular text
\begin{lemma}
blub; lemma is not endmarked
\end{lemma}
\begin{lemma}
abc
\begin{equation*}
abc
\end{equation*}
\begin{enumerate}
\item abc
\begin{itemize}
\item blub
\end{itemize}
\item jhdgd
\end{enumerate}
\end{lemma}
\end{theorem}

\begin{theorem}
regular text
\begin{remark}
blub; remark \emph{is} endmarked
\end{remark}
\end{theorem}


\begin{theorem}
regular text
\begin{remark}
blub; remark \emph{is} endmarked
\end{remark}
\begin{lemma}
abc; lemma is not endmarked
\begin{equation*}
abc
\end{equation*}
\begin{enumerate}
\item abc
\begin{itemize}
\item blub
\end{itemize}
\item jhdgd
\end{enumerate}
\end{lemma}
\end{theorem}




\clearpage

\subsection{proof tests}

\begin{proof}
regular text
\end{proof}

\begin{proof}
  This ends in a display
  \[
    3^2 + 4^2 = 5^2.
  \]
\end{proof}

\begin{proof}
  This ends in a display
  \[
    3^2 + 4^2 = 5^2.
  \]
more text
\begin{align}
a&=b \\
c&=d
\end{align}
%\[a=b\]
\end{proof}

\begin{proof}
  This ends in a display
  \[
    3^2 + 4^2 = 5^2.
  \]
more text
\begin{align*}
a=b
\end{align*}
abc
\end{proof}

\begin{proof}
  This ends with
  \begin{enumerate}
  \item an
  \item enumerated
  \item list.
  \end{enumerate}
\end{proof}

\begin{proof}
  This ends with
  \begin{enumerate}
  \item an
  \item enumerated
  \item list
  \item ending in equation:
  \begin{align*}
  a&=b
  \end{align*}
  \end{enumerate}
\end{proof}

\begin{proof}
  This ends with
  \begin{enumerate}
  \item an
  \item enumerated
  \item list.
  \begin{enumerate}
  \item abc
  \end{enumerate}
  \item blub
  \end{enumerate}
\end{proof}

\begin{proof}
  This ends with
  \begin{enumerate}
  \item an
  \item enumerated
  \item list.
  \begin{enumerate}
  \item abc
  \end{enumerate}
  \end{enumerate}
\end{proof}

\end{document}
